\documentclass[reqno, 12pt]{amsart}

\usepackage{amssymb}
\usepackage{amsthm}
\usepackage{amsmath}
\usepackage{amsxtra}
\usepackage{mathrsfs}
\usepackage{comment}
\usepackage{graphicx}
\usepackage{placeins}
\usepackage{multicol}
\usepackage{bbm}
\usepackage{stmaryrd}
%\usepackage{enumerate}
\usepackage[inline, shortlabels]{enumitem}
\usepackage{mathdots}
\usepackage{colonequals}
\usepackage{hyperref}


\usepackage{calc}
\usepackage{tikz}
\usetikzlibrary{arrows,calc,automata,shadows,backgrounds,positioning,intersections,fadings,decorations.pathreplacing,shapes,matrix}
\usepackage{tikz-cd}

\usepackage{fullpage}

\newtheorem{thm}{Theorem}[section]
\newtheorem{lem}[thm]{Lemma}
\newtheorem{cor}[thm]{Corollary}
\newtheorem{conj}[thm]{Conjecture}
\newtheorem{prop}[thm]{Proposition}
\theoremstyle{definition}  
\newtheorem{defn} [thm] {Definition} 
\newtheorem{claim}[thm]{Claim}
\newtheorem{example} [thm] {Example}
\newtheorem{rem} [thm] {Remark}

\newenvironment{problab}[1]
{\noindent\textbf{Problem #1}.}
{\vskip 6pt}

\newenvironment{exolab}[1]
{\noindent\textbf{Exercise #1}.}
{\vskip 6pt}

\theoremstyle{remark}
\newtheorem*{soln}{Solution}

\newcommand{\C}{\mathbb C}
\renewcommand{\H}{\mathbb H}
\newcommand{\D}{\mathbb D}
\newcommand{\F}{\mathbb F}
\newcommand{\Q}{\mathbb Q}
\newcommand{\R}{\mathbb R}
\newcommand{\Z}{\mathbb Z}
\newcommand{\T}{\mathbb T}
\renewcommand{\P}{\mathbb P}
\renewcommand{\O}{\mathcal O}
\newcommand{\m}{\mathfrak m}
\newcommand{\A}{\mathbb A}
\renewcommand{\SS}{\mathbb S}

\newcommand{\B}{\mathcal B}
\newcommand{\E}{\mathcal E}

\newcommand{\SSpan}{\text{span}}

\DeclareMathOperator{\lcm}{lcm}
\DeclareMathOperator{\img}{img}
\DeclareMathOperator{\Ann}{Ann}
\DeclareMathOperator{\Tor}{Tor}
\DeclareMathOperator{\tr}{tr}
\DeclareMathOperator{\Hom}{Hom}
\DeclareMathOperator{\End}{End}
\DeclareMathOperator{\Aut}{Aut}
\DeclareMathOperator{\Gal}{Gal}
\DeclareMathOperator{\sgn}{sgn}
\DeclareMathOperator{\ord}{ord}
\DeclareMathOperator{\ad}{ad}
\DeclareMathOperator{\coker}{coker}
\DeclareMathOperator{\rank}{rank}

\DeclareMathOperator{\id}{id}

\usepackage{mathpazo}

\everymath{\displaystyle}

\tikzset{commutative diagrams/.cd, arrow style = tikz, diagrams = {>=latex}}

\newenvironment{amatrix}[1]{%
  \left(\begin{array}{@{}*{#1}{c}|c@{}}
}{%
  \end{array}\right)
}

\usepackage{abstract}
\renewcommand{\abstractname}{}    % clear the title
\renewcommand{\absnamepos}{empty} % originally center
%\setlength{\abstitleskip}

\begin{document}
%\renewcommand{\abstractname}{\vspace{-\baselineskip}}

\title{18.700 Problem Set 2}
%\author{Évariste Galois}
\dedicatory{Due Wednesday, September 18 at 11:59 pm on \href{https://canvas.mit.edu/courses/27315}{Canvas}}
%\thanks{}
%\contrib{}
%\address{}

\maketitle

\begin{abstract}
  \noindent Collaborated with:\\
  Sources used: 
\end{abstract}
\vskip 0.5cm

\noindent
\textbf{Note:} When using one or more of the vector space axioms in your solutions, make sure to indicate which axioms are being used and how. E.g., ``\dots where the second equality follows by associativity of scalar multiplication.''
\vskip 0.5cm

\begin{problab}{1} (6 points)
  \begin{enumerate}[(a)]
    \item 
      Suppose a $3 \times 5$ \emph{coefficient} matrix for a linear system has $3$ pivot columns. Is the system consistent? Why or why not?
    \item
      Suppose a linear system has a $3 \times 5$ \emph{augmented} matrix whose fifth columns is a pivot column. Is the system consistent? Why or why not?
    \item
      Suppose the coefficient matrix of a system of linear equations has a pivot position in every row. Explain why the system must be consistent.
  \end{enumerate}
\end{problab}

\begin{problab}{2} (1B, \#2) (5 points)
  Suppose $a \in \F, v \in V$, and $av = 0$. Prove that $a=0$ or $v=0$.
\end{problab}

\begin{problab}{3} (4 points)
  Let $W$ be the union of the first and third quadrants in $\R^2$, i.e.,
  \begin{equation*}
    W \colonequals \{(x,y) \in \R^2 : xy \geq 0\} \, .
  \end{equation*}
  \begin{enumerate}[(a)]
    \item Show that $cw \in W$ for all $c \in \R$ and all $w \in W$.
    \item Is $W$ a vector space? Why or why not?
  \end{enumerate}
\end{problab}

\begin{problab}{4} (1B, \#6) (3 points)
  Let $\infty$ and $-\infty$ denote two distinct objects, neither of which is in $\R$. Define an addition and scalar multiplication on $\R \cup\{\infty,-\infty\}$ as follows. The sum and product of two real numbers is as usual, and for $t \in \R$ define
$$
t \infty=
\begin{cases}
  -\infty & \text { if } t<0, \\
  0 & \text { if } t=0, \\
  \infty & \text { if } t>0,
\end{cases}
\qquad
t(-\infty)=
\begin{cases}
  \infty & \text { if } t<0, \\
  0 & \text { if } t=0, \\
  -\infty & \text { if } t>0,
\end{cases}
$$
and
$$
\begin{aligned}
  t+\infty & =\infty+t=\infty+\infty=\infty, \\
  t+(-\infty) & =(-\infty)+t=(-\infty)+(-\infty)=-\infty, \\
  \infty+(-\infty) & =(-\infty)+\infty=0 .
\end{aligned}
$$
With these operations of addition and scalar multiplication, is $\R \cup\{\infty,-\infty\}$ a vector space over $\R$? Explain.
\end{problab}

\begin{problab}{5} (6 points)
  Let $U = \{0\}$ and define $0 + 0 = 0$ and $\lambda \, 0 = 0$ for all $\lambda \in \F$. Prove that $U$ is a vector space by verifying each of the vector space axioms.
\end{problab}


\end{document} 	
